\section{Clearing denominators}\label{sec:arithmetic}

\subsection{The progression estimate}

We first clear the coefficients of a single numerator block.
This supplies the divisibility needed when we expand at each pole.
Put $d_n=\lcm(1,\ldots,n)$ and $\delta_n=2d_n$.
For $\theta=v/d\in\Theta$ in lowest terms, consider
\[
 B_\theta(t)=\frac{A_d^n}{(n!)^L}
       \prod_{j=0}^{Ln-1}(t-2rn+2j+v/d).
\]

\begin{lemma}\label{lem:progression}
For every integer $v_0$ and every $\ell\ge0$,
\[
 d_n^\ell[t^\ell]B_\theta(t+v_0)\in\Z.
\]
\end{lemma}
\begin{proof}
Clear $d$ from each factor using $d^{Ln}$ in $A_d^n$.
At an odd prime $q\nmid d$, the constants form a progression
with step $2d$, a unit modulo $q$. Split the $Ln$ factors into
$L$ blocks of $n$. At each level $q^e\le n$, at least
$L\lfloor n/q^e\rfloor$ constants are divisible.
Extracting degree $\ell$ replaces $\ell$ constants by variable
terms, losing at most $\ell$ contributions at that level.
The valuation is at least
\[
 L v_q(n!)-\ell v_q(d_n).
\]
Levels above $n$ can only improve the bound. A term retaining
a zero constant is zero and needs no estimate.

If $q\mid d$, the explicit factor
$q^{n\lceil L/(q-1)\rceil}$ pays for $L v_q(n!)$.
At $q=2$, the polynomial after clearing $d$ has integer
coefficients, and the factor $2^{Ln}$ pays for $L v_2(n!)$.
These bounds prove the claim.
\end{proof}
Products and integer powers preserve this coefficient rule
by convolution.

\subsection{The coefficient bounds}

\begin{proposition}\label{prop:denominators}
For every positive odd $n$,
\begin{align}
 \delta_n^{a-i}c_{i,h}(n)&\in\Z,\label{eq:c-integral}\\
 \delta_n^{a+k-\nu}A_{\nu,h}(n),\
 \delta_n^{a+k-\nu}\rho_\nu(n)&\in\Z,\label{eq:rho-integral}\\
 \delta_n^{a+k}q_b(n)&\in\Z
                       \qquad(b=1,3,\ldots,M).\label{eq:q-integral}
\end{align}
\end{proposition}
\begin{proof}
Write $B_n=K^{mn}P_n^m/(n!)^{mDL}$ and suppress the index $n$
on the coefficients. Expanding at $-2h$ gives
\begin{equation}\label{eq:local-coefficient}
 \begin{split}
 c_{i,h}=\binom nh^a[t^{a-i}]
 \left\{W_n(t-2h)B_n(t-2h)
       \prod_{\ell\ne h}\left(1+\frac{t}{2(\ell-h)}\right)^{-a}
 \right\}.
 \end{split}
\end{equation}
The constant factors from the other poles have product
$2^{an}(h!(n-h)!)^a$, since $a$ is even.
Their cancellation with $\kappa_n$ gives the displayed prefactor.
Each $2(\ell-h)$ divides $\delta_n$. The progression lemma,
the integer coefficients of $W_n$, and the binomial series
prove \eqref{eq:c-integral}.
Formula~\eqref{eq:partials} then gives \eqref{eq:rho-integral},
because $a+k-(i+k)=a-i$.

The finite sums defining $q_b$ may contain prime powers not
covered by $\delta_n$, so individual summands need not be integral
after scaling. We use the surviving zeros to clear them in groups,
as in~\cite[Section~4]{Tcount}.
Fix an odd $j<2bh$ and put
\[
 t_0=j/b-2h,\qquad x_\ell=j/b+2(\ell-h).
\]
The point $t_0$ is a non-pole grid zero in $(-2n,0)$.
Thus $V_n(t_0)=0$ gives
\begin{equation}\label{eq:grouping}
 \sum_\nu\frac{A_{\nu,h}}{x_h^\nu}
       =-\sum_{\ell\ne h}\sum_\nu\frac{A_{\nu,\ell}}{x_\ell^\nu}.
\end{equation}
At a prime $q$ put $v=v_q(\delta_n)$. At most one $x_\ell$
has valuation greater than $v$: the difference of two such
numbers would be the nonzero integer $2(\ell-\ell')$,
which divides $\delta_n$. This argument also applies when
$q$ divides the grid denominator $b$.
From \eqref{eq:rho-integral},
\[
 v_q\!\left(\frac{\delta_n^{a+k}A_{\nu,\ell}}{x_\ell^\nu}\right)
                 \ge\nu\bigl(v-v_q(x_\ell)\bigr).
\]
If $v_q(x_h)\le v$, the left group in \eqref{eq:grouping}
is integral after scaling. Otherwise every group on the right
is integral after scaling, so the left group is too.
Apply this at every prime, and sum the groups with their signs
in \eqref{eq:q}. This proves \eqref{eq:q-integral}.
\end{proof}

The common multiplier is $\delta_n^{a+k}$.
Differentiation shifts an order-$i$ coefficient to order $i+k$,
while its denominator bound remains $\delta_n^{a-i}$.
This explains the additional exponent $k$.
